\chapter{Before traveling}

In this section you can expect all the things you should
do before you depart. This includes hardening some devices
and making sure you have a backup of the data you'll be
traveling with. It also covers letting someone know you'll
be calling in case of an emergency and how to identify
you.

If you're already reading this while traveling, don't
worry, this will still be helpful, it just might be a bit
more challenging depending on your context.

To start, you can remove or securely store any data,
devices, printed files, and documents you don't absolutely
need on the trip. The less sensitive information and fewer
critical assets you carry, the lower the risk and impact
if loss, theft, or inspection occurs. Minimize your
digital and physical footprint by leaving backups and
originals securely at home or in trusted locations, and
encrypt what must travel with you.

This principle applies to laptops, phones, hardware
wallets, paper backups, and any portable storage media.

\section{Perform a quick threat model}

Take five minutes to map out your risks. Identify what
assets you're carrying (laptops, phones, hardware wallets,
seed phrases, account access), who might target your
assets (thieves, cybercriminals, border agents, etc.).
Lastly, identify how attacks might happen (device theft,
tampering, malware, coercion). For each threat, plan a
mitigation\footnote{\url{https://youtu.be/igNa7z4hwMQ?t=2465}}
if you feel this threat is applicable to you. For example,
if you're carrying a hardware wallet, a threat is
pick-pocketing -- mitigation could be keeping it attached
to yourself (just don't use Ledger's necklace visibly) and
securing it with PIN/passphrase (no patterns, not repeated
across devices).

This exercise keeps security measures proportionate to
your situation.

\section{Secure devices with encryption \& updates}

Enable full-disk encryption on all devices (laptops,
phones, tablets) to protect data if lost or stolen. Most
modern OSes have them installed, not always enabled by
default (e.g.\ BitLocker for Windows, FileVault for
MacOS, LUKS for Linux, or Android/iOS encryption -- just
ensure a strong password/PIN is set).

Use popular devices like iPhones and Pixels. Install the
latest OS and app updates since these patch security
vulnerabilities.\footnote{A recent attack named Pixnapping
has been uncovered, affecting many Android devices,
particularly Pixels -- \url{https://www.pixnapping.com}}
Latest is not always safest, but security patches usually
win that trade-off. You can also consider to install --
only trustworthy -- mobile security
applications,\footnote{On iOS, due to sandboxing, apps
can't scan the entire OS for malware, so many security
apps focus on VPN, phishing web protection, breach
alerts, etc.} which try to add a layer of control and also
remind you of important security updates.

If you must bring high-risk confidential data, consider
encrypting those files individually using tools like
VeraCrypt, FileVault, BitLocker, LUKS/dm-crypt, or others.

\section{Back up and prepare for loss}

Back up your devices before travel (and ensure cloud
backups like iCloud/Google are up to date). This way, if a
device is lost or confiscated, you won't lose important
information.

Record device details (make, model, serial numbers, IMEI
for phones) and keep that list separate -- it will help in
filing reports or remote-wipe commands.

If your company uses Mobile Device Management (MDM) on
phones or laptops, verify the device is enrolled and you
know how to trigger a remote wipe or ``lost mode.'' Test
``Find My'' or equivalent device-finding services so you
can use them if needed. Pack chargers and cables so you
won't need to borrow unknown ones (which could be
malicious).

\section{Protect Accounts with 2FA redundancy}

Plan for how you'll access accounts that use two-factor
authentication if your main device is unavailable. For
authenticator apps (TOTP codes), export backup codes for
your important accounts and keep them securely in services
like 1Password or NordPass (you can check
passwordmanager.com for more information on which one to
pick).

Alternatively, consider storing them elsewhere
physically.\footnote{You can write them on paper and keep
them in a safe, or store them on an encrypted drive.} If
your 2FA is tied to a phone number (SMS or voice), disable
SMS for 2FA. For technologies that are unfortunately
dependent on phone numbers, use a separate line (not your
personal one) from services like Google Voice, Burner App
or SLYFONE. Transfer your number to an eSIM since they are
harder to physically steal or swap than a physical
SIM.\footnote{\url{https://support.apple.com/en-ca/118227}}

Ensure your password manager is accessible -- many like
1Password have a Travel Mode that removes sensitive vaults
from your device while you travel (you can restore them
later). This limits exposure if your device is searched or
seized.\footnote{\url{https://www.schneier.com/blog/archives/2025/04/cell-phone-opsec-for-border-crossings.html}}

You can also register a backup or alternative 2FA method,
like a secondary phone (the device) or a PIN-protected
hardware key to be used as a passkey. Be responsible with
the use of software passkeys, particularly in situations
where you are using a password manager to store both your
passwords and passkeys, as this creates a single point of
failure.

\section{Secure your wallets and keys}

\subsection{Hardware wallets (e.g.\ Ledger, Trezor)}

Update the firmware and test the device before you leave,
don't do this while traveling. Do NOT bring any written
seed phrases under any circumstance -- seed backups are
unencrypted keys that are far easier to steal or copy than
a hardware device. Leave seed backups in a secure place,
travel only with your hardware wallet if necessary. Enable
all security features on the device (set a strong PIN, and
use a BIP39 passphrase for example, if supported) so that
even if the device is stolen, the amount of information
required to access your crypto is high. Carry hardware
wallets and security keys in your carry-on or under your
sight, not in checked luggage, to avoid loss or tampering.

\subsection{Yubikeys and 2FA tokens}

Bring them to protect logins and make sure they're enabled
on your critical accounts. Keep them on your person or in
a separate bag from your laptop/phone so that a thief or
snoop can't easily steal both at once. If you have a spare
hardware wallet or Yubikey, you might leave one at home as
a backup in case the one you carry is lost. Add, when
possible, an extra layer of security to the token, such as
a PIN code.

\section{Lock down phones}

On your smartphone, take advantage of security settings
before you travel. Set a strong passcode (not just a
4-digit PIN or pattern). Consider disabling Touch ID/Face
ID if you might face situations where someone could
force-unlock your device using your biometrics -- in many
jurisdictions, authorities can compel a fingerprint or
face scan more easily than a memorized
password.\footnote{\url{https://www.schneier.com/blog/archives/2025/04/cell-phone-opsec-for-border-crossings.html}}
At a minimum, know how to quickly and temporarily disable
biometrics;\footnote{\url{https://news.ycombinator.com/item?id=43630624}}
for example, on an iPhone, holding the side button and a
volume button will trigger an emergency mode that requires
the passcode to
unlock.\footnote{\url{https://www.themacguys.com/essential-travel-security-tips-for-apple-devices}}
On Android, use Lockdown mode if available (which only
temporarily disables biometrics and is different from iOS
Lockdown Mode).

If you have an iPhone, enable Lockdown Mode (extreme
protection on iOS) even if you are not at high risk or
traveling to a high-threat region -- but be aware it
restricts many
features,\footnote{\url{https://support.apple.com/en-gb/105120}}
though totally worth it.

Disabling or restricting USB debugging on Android and
using iOS USB restricted mode helps prevent unauthorized
physical USB access and reduces risks from malicious
cables or compromised charging stations.

If your phone is managed by work (MDM), inform your IT
team of your travel so they can assist with any
location-based security policies and ensure you have the
latest security profile. Finally, consider using a
dedicated eSIM/local SIM for travel data. This can protect
your primary phone number (you can keep your main line on
eSIM and turn it off, while using a local data eSIM for
mobile internet) and avoids physical SIM issues.

\subsection{Configure additional phone protections}

On iOS devices, disable Control Center and Notification
Center access from the lock screen (Settings $>$ Face ID
\& Passcode) to prevent thieves from seeing notifications
or enabling Airplane Mode without unlocking. Disable USB
accessory connections when locked to prevent unauthorized
connections. For device recovery, ensure Find My iPhone is
enabled with ``Send Last Location'' and ``Find Network''
options activated so tracking continues even if the device
is powered off.

On Android, similar protections can be configured: disable
notifications on lock screen (Settings $>$ Notifications
$>$ On lock screen $>$ Don't show notifications), and
enable ``Find My Device'' with all location services.

\subsection{PIN security}

Many financial apps default to PIN verification instead of
biometrics, which means a thief who has your phone and
knows your PIN could potentially access your financial
accounts even if they can't bypass Face ID/Touch ID. Use
unique PINs for banking apps that differ from your device
PIN, or where possible, configure these apps to require
biometric verification for every login.

\section{Minimize digital \& social media footprint}

It's not just cyber threats -- operational security
matters. Avoid announcing your travel plans publicly on
social
media\footnote{\url{https://blackcloak.io/social-media-and-travel-be-careful-of-what-you-share/}}
and be careful with real-time updates. Posts about being
away from home can signal to criminals that you (and
possibly valuable devices or even your house) are
vulnerable. Consider sharing trip photos and crypto
conference highlights after you return or only with
trusted contacts. If you post at all, delay the location
by a day or two. Ensure your social media privacy settings
are tightened so strangers can't see your travel posts.

When registering for events, use privacy-focused tools
like iOS's Hide My Email or create burner emails through
providers like Tuta. ProtonMail has had some controversy
lately, but could still be a good alternative. If you
don't need to keep the inbox at all, just create a throw
away email account, there are plenty of options out there.

Avoid giving out personal details unnecessarily during
registration -- use minimal, generic info to reduce your
digital footprint.

Discretion is key: if possible, don't advertise that you
work in crypto\footnote{\url{https://www.unchained.com/blog/7-tips-traveling-bitcoin}}
or carry cryptocurrency. For example, remove or cover any
crypto stickers on laptops or bags, and avoid wearing
company swag or Bitcoin/Ethereum logos while in transit.
These can be neon signs attracting thieves or unwanted
attention (``I have valuable data or wallets!''). If asked
about your work or luggage by curious strangers, have a
cover story (e.g.\ ``I work in finance/IT'') rather than
``I manage a crypto fund''. High-profile team members
might travel under pseudonyms or at least not list their
company on luggage tags to stay low-profile. Also, be wary
that you might be traveling with someone with these
characteristics, so don't give them away.

\section{Plan emergency and incident responses}

Before departure, know what you'll do if something goes
wrong. Have a fallback plan in case of device loss: who
will you notify \& how (have safe words and beware of
deepfakes or impersonations); how will you revoke access
to sensitive accounts; and how will you continue working
(e.g., a backup laptop or a colleague who can step in).

If traveling to countries with strict
tech\footnote{\url{https://www.perplexity.ai/search/countries-that-list-cryptopgra-bBSItefPSUi7HP2lTNekjA}}
or encryption
laws\footnote{\url{https://it.cornell.edu/security-and-policy/travel-internationally-technology}}
(e.g., China, Russia, UAE), devices like VPNs, encrypted
messaging apps (Signal, or Telegram with Secret Chats
only), hardware wallets (Ledger, Trezor), Yubikeys, or
encryption software (VeraCrypt, BitLocker) may be flagged
by border authorities. Research local laws beforehand.
Consider carrying a travel letter from your organization
explaining the professional need for these tools, or use
sanitized loaner devices to avoid issues at border
controls.

Share your itinerary and contact information with a
trusted peer so they can assist or monitor for any issues.

Finally, schedule critical work (especially high-value
transactions) for before or after your trip if possible,
so you're not forced to do ultra-sensitive actions on the
road. Criminals usually play with the ``time-sensitive
factor,'' trying to trick you into doing something quick
and urgent, by committing something reckless.

There's more information for high profile individuals in
the last chapter, page~\pageref{ch:high-profile}.
